# Architecture Decision Record (ADR-003)

## Наименование
Перевод вычислительного ядра с текстовых строк на сырые бинарные буферы (`node:buffer`).

## Статус
Утверждено.

## 1. Контекст и Постановка Задачи
В исходном прототипе конвейера обработка экранированных полей CSV выполнялась на уровне текстовых строк JavaScript с использованием регулярных выражений (Regular Expressions) и метода `.replace()`. 

В ходе тестирования выявлены три системных дефекта:
1. **Накладные расходы памяти (ConsStrings):** Специфика движка Google V8 при операциях `.slice()`, `.replace()` и конкатенации строк приводит к созданию виртуальных деревьев ссылок в куче, вызывая массовые кэш-промахи ЦП (Cache Misses) и нагрузку на сборщик мусора (Garbage Collector).
2. **Паразитная сложность $\mathcal{O}(K^2)$:** Прогон масок регулярных выражений по каждому входящему байту плоского текста приводил к холостым тактам ЦП на тривиальных данных.
3. **Индексный сдвиг:** Попытка удерживать глобальный указатель файла одновременно с локальными сдвигами регулярных выражений вызывала потерю координат и зацикливание автомата.

## 2. Принятое Архитектурное Решение
Полный отказ от использования строкового примитива внутри парсера. Перевод конвейера на сырые бинарные буферы (`node:buffer`) и внедрение алгоритма блочно-оконного сканирования (`Window Parsing`) [RFC 4180].

Реализована концепция полной автономии марковского ядра:
* Линейный кассетный автомат Вирта выступает ведущим драйвером потока, сканируя ASCII-коды.
* При обнаружении открывающей кавычки поля (`0x22`) автомат Вирта передает управление Маркову.
* Марков автономно выкачивает порцию данных в изолированный буфер `tmp`, находит легальный финал ячейки (`",` или `"\n`) и возвращает только примитивный указатель `new_index` [RFC 4180].

## 3. Математическое Обоснование (Свертка Каскада)
Выход на уровень бинарных данных позволил установить, что первая внешняя кавычка отсекается автоматом Вирта на входе, а все внутренние кавычки текста ячейки всегда идут строго парами `[0x22, 0x22]` [RFC 4180].

Данный инвариант позволил ликвидировать громоздкие многоэтажные маркеры состояний и свернуть весь каскад Нормальных Алгоритмов Маркова (НАМ) до двух правил строго по два байта:
1. Найти внутреннюю легальную пару `[0x22, 0x22]` $\rightarrow$ заменить на однобайтовый маркер `[0x01]`.
2. Найти финальный терминатор поля перед запятой $\rightarrow$ выдать сигнал `THE_END`.

## 4. Последствия и Эффективность
1. **Zero-Overhead по памяти:** Так как длина левой части марковского правила ($2\text{ байта}$) стала в точности равна длине правой части, размер буфера `tmp` не меняется во время мутаций. Из кода полностью исключаются вызовы `Buffer.concat` [RFC 4180].
2. **Локальность кэша:** Замена выполняется in-place через метод `rule.copy(tmp, offset)` прямо в один такт процессора. ЦП крутит чистые ассемблерные инструкции сравнения регистров `CMP` на локальном участке L1-кэша.
3. **Безопасность рантайма:** Выделение памяти в куче V8 сведено к нулю. Сборщик мусора полностью исключен из процесса вычислений.
